\documentclass[11pt, a4paper]{article}

\usepackage[a4paper, top=2.5cm, bottom=2.5cm, left=2cm, right=2cm]{geometry}
\usepackage{amsmath}
\usepackage{amssymb}
\usepackage[colorlinks=true, linkcolor=blue, urlcolor=blue, citecolor=blue]{hyperref}

\title{Theory of Everything - Volume 46 \\ \Large The Simulation Hypothesis}
\author{Omega Protocol}
\date{\today}

\begin{document}

\maketitle

\section*{Layman's Summary}
Are we living in a computer simulation? Volume 46 proves that the question itself is meaningless. Physics is computation. A perfectly simulated universe is mathematically identical to a "real" universe. However, we also prove that if we *are* inside a simulation running on a finite computer, we should be able to detect the glitches.

\section{Computational Equivalence (Axiom 49)}
We establish that there is no difference between "base reality" and a perfect simulation. Both are simply mathematical manifolds of informational flow. If an intelligence inside a computer experiences relative entropy dissipation, their reality is as mathematically "real" as the universe hosting the computer.

\section{The Simulation Bound (Theorem)}
You cannot perfectly simulate a computer inside a computer of the exact same size. We prove that any embedded simulation must operate at a lower resolution (a lower Quantum Fisher Information Metric) than the universe hosting it. This strict thermodynamic bound prevents an infinite Russian doll scenario of simulations within simulations.

\section{Glitches and Anomalies (Theorem)}
If the computer running our universe has finite RAM and processing power, it has to cut corners. We prove that these computational shortcuts—truncation errors—will inevitably manifest in our universe as high-energy quantum anomalies. By looking closely at the behavior of ultra-high-energy cosmic rays, we can theoretically test if our reality is being "rendered" on a finite machine.

\end{document}
